\section{Introduction}\label{sec:intro}

AI agents now write, test and repair code in loops with little human involvement. The
artifacts they manipulate were designed for human authors: a program is a set of text files
whose identity is a path, whose meaning depends on names and declaration order, whose
diagnostics are prose, and whose compiled behaviour may depend on undefined behaviour and on
the floating-point unit of the target. Verification-aware languages such as
Dafny~\cite{leino2010dafny}, Why3~\cite{filliatre2013why3}, F*~\cite{swamy2016fstar} and
Verus~\cite{lattuada2023verus} close part of the gap: a program carries a contract that a
solver checks, and language models are becoming competent at writing such
programs~\cite{bursuc2025vericoding}. Identity, error reporting, determinism and composition
nevertheless remain organized around text written for people.

\VG{} starts from a different question: what should a program representation look like if
every design decision is judged by what it does for an agent that writes, checks and repairs
code with no human in the loop? Its answer has four parts.
\begin{enumerate}
\item A program is a set of \emph{nodes}. Each node is one function: a typed,
  single-assignment dataflow graph of combinators, with a contract (pre- and postconditions,
  loop invariants, cost bounds) and an effect set. Its identity is a BLAKE3
  hash~\cite{blake3} of its canonical structure and contract, so equal programs have equal
  names on every machine and every proof can be cached by content.
\item Agents write a textual projection, VG-Sym, in which they name things with labels; the
  toolchain computes every hash. Ports and labels are defined before use, so every prefix of
  a VG-Sym file can be checked.
\item Each contract is split into \emph{obligations}, each discharged by one quantifier-free
  SMT query in a fragment that Z3~\cite{demoura2008z3}, cvc5~\cite{barbosa2022cvc5} and
  Bitwuzla~\cite{niemetz2023bitwuzla} all accept. A failure comes back as structured data: an
  obligation identifier, a counterexample that a reference interpreter replays, and a failing
  slice of the graph.
\item Verified nodes are meant to compile, through textual LLVM IR~\cite{lattner2004llvm}, to
  x86\_64, ARM64 and WebAssembly binaries that agree bit for bit with the interpreter, with
  proven facts turned into optimization hints and an operating-system sandbox derived from the
  capabilities of the root node.
\end{enumerate}

These choices raise precise questions that an informal specification answers only loosely.
When do two programs share a hash, and when must they not? Can effectful operations be merged
by hash-consing? What does a floating-point value \emph{mean} when x86\_64, AArch64 and
WebAssembly disagree on the bits of NaN? Which proof rules make bounded loops verifiable with
quantifier-free queries, and are they sound? What does a counterexample that does not replay
tell the agent? This paper answers these questions with definitions, theorems and proofs.

\paragraph{Contributions.}
\begin{itemize}
\item A core calculus for \VG{} nodes (\cref{sec:calculus}): types with a one-NaN float
  domain, nodes as term graphs with linear effect tokens, a strict evaluation semantics with
  traps, contracts with explicitly unspecified values, and two confinement results: a node's
  effects stay within its declared set, and capabilities cannot be forged.
\item A characterization of content addressing (\cref{sec:identity}). Under an explicit
  field-capacity condition and the usual collision-freeness idealization, two bodies have the
  same hash if and only if their unravelings coincide up to commutativity
  (\cref{thm:hash-char}); equal node hashes imply equal meaning (\cref{thm:hash-meaning}); and
  hash-consing never merges two effectful operations when effect tokens are linear
  (\cref{lem:eff-never-merged}). We also prove that no polynomial-time canonical form
  identifies exactly the semantically equivalent nodes unless $\mathrm{P}=\mathrm{NP}$
  (\cref{thm:conp}), which justifies the design's identity by specification.
\item A one-NaN semantics for IEEE~754 (\cref{sec:float}). Identifying all NaN bit patterns
  makes every float operation of the minimal viable product (MVP) a congruence, except
  copySign with a NaN sign operand, which the design therefore treats as partial
  (\cref{thm:float-congruence}). Under this semantics addition, multiplication, minimum,
  maximum and the product part of a fused multiply-add are commutative bit for bit
  (\cref{thm:float-comm}), which licenses sorting their operands during canonicalization.
\item Soundness of modular verification (\cref{sec:verification}). The combinator proof
  rules, the fresh-element rule and an exact bit-width encoding of mathematical integers
  (\cref{thm:widths}) are sound (\cref{thm:soundness}), provided each obligation is checked
  only against facts of operations that precede it in a topological order. A four-operation
  program shows that this proviso is necessary (\cref{prop:ordering-needed}); the current
  specification does not state it. The counterexample-replay classification is sound too
  (\cref{thm:replay}).
\item A complete IEEE~754 proof that the case study's soft-knee kernel meets its contract,
  and an exact characterization of the inputs on which its first draft fails
  (\cref{sec:case}), cross-checked exhaustively on every finite binary32 input for ten knee
  values, $4.28\times10^{10}$ evaluations.
\item An honest account of the prototype (\cref{sec:status}): what exists, what it passes, and
  the findings of an independent review, two of which we reproduce and show to be exactly
  violations of the side conditions of our theorems.
\end{itemize}

\paragraph{Status and scope.}
\VG{} is at an early stage. Its core intermediate representation (parser, printer,
canonicalizer, hashing, binary codec and object store) is implemented in about 9,600 lines of
Rust, tests included. The type checker, interpreter, verifier, code generator, agent loop,
registry and lifter are specified but not implemented. Our theorems are about the
specification and the core calculus, not about running code; where we report a measurement,
we name the commit and the command. We make no claim about verification time, the speed of
generated code or synthesis success rates, because none of them has been measured.

\paragraph{Organization.}
\Cref{sec:overview} introduces \VG{} through its case study. \Cref{sec:calculus} defines the
core calculus. \Cref{sec:identity} studies identity and content addressing,
\cref{sec:float} floating point, and \cref{sec:verification} verification.
\Cref{sec:case} proves the case study correct. \Cref{sec:status} reports the prototype,
\cref{sec:related} discusses related work and \cref{sec:discussion} the limitations.
